\section{Free fundamental graded Lie algebras}\label{section2}

First of all we give several def\/initions about graded Lie algebras.
Let $\mathfrak g$ be a Lie algebra. Assume that there is given a family of
subspaces $(\mathfrak g_p)_{p\in\mathbb Z}$ of $\mathfrak g$ satisfying
the following conditions:
\begin{enumerate}\itemsep=0pt
\renewcommand{\labelenumi}{$(\roman{enumi})$}
\item $\gla g$;
\item $\dim \mathfrak g_p<\infty$ for all $p\in\mathbb Z$;
\item $[\mathfrak g_p,\mathfrak g_q]\subset \mathfrak g_{p+q}$
for all $p,q\in\mathbb Z$.
\end{enumerate}
Under these conditions, we say that $\gla g$ is a graded Lie algebra (GLA).
Moreover we def\/ine the notion of homomorphism, isomorphism, monomorphism,
epimorphism, subalgebra and
ideal for GLAs in an obvious manner.

A GLA $\gla g$ is called transitive if
for $X\in\mathfrak g_p$ $(p\geqq0)$, $[X,\mathfrak g_-]=\{0\}$ implies $X=0$,
where $\mathfrak g_-$ is the negative part
$\bigoplus\limits_{p<0}\mathfrak g_p$ of $\mathfrak g$.
Furthermore a GLA $\gla g$ is called irreducible if
the $\mathfrak g_0$-module $\mathfrak g_{-1}$ is irreducible.

Let $\mu$ be a positive integer. A GLA
$\gla g$ is said to be of depth $\mu$ if $\mathfrak g_{-\mu}\ne\{0\}$ and
$\mathfrak g_p=\{0\}$ for all $p<-\mu$.

Next we def\/ine fundamental GLAs.
A GLA $\mathfrak m=\bigoplus\limits_{p<0}\mathfrak g_p$ is
called a \textit{fundamental graded Lie algebra} (FGLA)
if the following conditions hold:
\begin{enumerate}\itemsep=0pt
\renewcommand{\labelenumi}{$(\roman{enumi})$}
\item $\dim \mathfrak m<\infty$;
\item $\mathfrak g_{-1}\ne\{0\}$, and $\mathfrak m$ is generated by $\mathfrak g_{-1}$,
or more precisely $\mathfrak g_{p-1}=[\mathfrak g_p,\mathfrak g_{-1}]$
for all $p<0$.
\end{enumerate}
If an FGLA $\mathfrak m=\bigoplus\limits_{p<0}\mathfrak g_p$
is of depth $\mu$, then $\mathfrak m$
is also said to be of the $\mu$-th kind.
Moreover an FGLA $\mathfrak m=\bigoplus\limits_{p<0}\mathfrak g_p$ is called non-degenerate if
for $X\in\mathfrak g_{-1}$, $[X,\mathfrak g_{-1}]=\{0\}$ implies $X=0$.
\par
Let $\mathfrak m=\bigoplus\limits_{p<0}\mathfrak g_p$ be an FGLA of
the $\mu$-th kind, where $\mu\geqq2$. $\mathfrak m$ is called a free
fundamental graded Lie algebra of type $(n,\mu)$
if the following conditions hold:
\begin{enumerate}\itemsep=0pt
\renewcommand{\labelenumi}{$(\roman{enumi})$}
\item $\dim \mathfrak g_{-1}=n$;
\item Let $\mathfrak m'=\bigoplus\limits_{p<0}\mathfrak g'_p$ be an FGLA
of the $\mu$-th kind and let $\varphi$ be a surjective linear
mapping of $\mathfrak g_{-1}$ onto $\mathfrak g'_{-1}$.
Then $\varphi$ can be extended uniquely to a GLA
epimorphism of $\mathfrak m$ onto $\mathfrak m'$.
\end{enumerate}

\begin{proposition}\label{prop2.1} Let $n$ and $\mu$ be positive integers such that $n,\mu\geqq2$.
\begin{enumerate}\itemsep=0pt
\item There exists a unique free FGLA of type
$(n,\mu)$ up to isomorphism.
\item Let $\mathfrak m=\bigoplus\limits_{p<0}\mathfrak g_p$ be a free
FGLA of type $(n,\mu)$.
We denote by $\Der(\mathfrak m)_0$ the Lie algebra of all the derivations
of $\mathfrak m$ preserving the gradation of $\mathfrak m$.
Then the mapping $\Phi:\Der(\mathfrak m)_0\ni D\mapsto D|\mathfrak g_{-1}\in\gl(\mathfrak g_{-1})$ is a Lie algebra isomorphism.
\end{enumerate}
\end{proposition}
\begin{proof} (1)
The uniqueness of a free FGLA of type $(n,\mu)$ follows from the def\/inition.
We set $X=\{1,\dots,n\}$. Let $L(X)$ be the free Lie algebra on $X$
(see \cite[Chapter~II, \S~2]{bou72:1})
and let $i:X\to L(X)$ be the canonical injection.
We def\/ine a mapping $\phi$ of $X$ into $\mathbb Z$ by $\phi(k)=-1$ $(k\in X)$.
The mapping $\phi$ def\/ines the natural gradation $(L(X)_p)_{p<0}$
on $L(X)$ such that: $(i)$ $L(X)$ is generated by $L(X)_{-1}$;
$(ii)$ $\{i(1),\dots,i(n)\}$ is a basis of $L(X)_{-1}$
(see \cite[Chapter~II, \S~2, no.~6]{bou72:1}). Note that if $n>1$, then
$L(X)_p\ne0$ for all $p<0$.
We set $\mathfrak a=\bigoplus\limits_{p<-\mu}L(X)_p$;
then $\mathfrak a$ is a graded ideal of $L(X)$ and the factor GLA $\mathfrak m=L(X)/\mathfrak a$ becomes an FGLA of the $\mu$-th kind.
We put $\mathfrak a_p=\mathfrak a\cap L(X)_p$ and
$\mathfrak g_p=L(X)_p/\mathfrak a_p$.

Now we prove that $\mathfrak m=\bigoplus\limits_{p<0}\mathfrak g_p$ is a free
FGLA of type $(n,\mu)$.
Let $\mathfrak m'=\bigoplus\limits_{p<0}\mathfrak g'_p$ be an FGLA of the
$\mu$-th kind and let
$\varphi$ be a surjective linear mapping of $\mathfrak g_{-1}$ onto
$\mathfrak g'_{-1}$.
Let $h$ be a mapping of~$X$ into $\mathfrak m'$ def\/ined by $h(k)=\varphi(i(k))$ $(k\in X)$.
Then there exists a Lie algebra homomorphism~$\tilde{h}$ of~$L(X)$
into $\mathfrak m'$ such that $\tilde{h}\circ i=h$.
Since $L(X)$ (resp.~$\mathfrak m'$) is generated by
$L(X)_{-1}$ (resp.~$\mathfrak g'_{-1}$), $\tilde{h}$ is surjective.
Since $\mathfrak m'=\bigoplus\limits_{p<0}\mathfrak g'_p$ is of the
$\mu$-th kind, $\tilde{h}(\mathfrak a)=0$, so
$\tilde{h}$ induces a GLA epimorphism~$L(\varphi)$ of
$\mathfrak m$ onto~$\mathfrak m'$ such that
$L(\varphi)|\mathfrak g_{-1}=\varphi$.
The homomorphism~$L(\varphi)$ is unique, because
$\mathfrak m=\bigoplus\limits_{p<0}\mathfrak g_p$ is generated by
$\mathfrak g_{-1}$.
Thus $\mathfrak m$ is a free FGLA of type~$(n,\mu)$.

(2) Assume that $\mathfrak m$ is a free FGLA
constructed in~(1).
Let $\phi$ be an endomorphism of $\mathfrak g_{-1}$.
By Corollary to Proposition~8 of \cite[Chapter~II, \S~2, no.~8]{bou72:1},
$\phi$ can be extended uniquely to a~unique derivation $D$ of $L(X)$.
Since $D(L(X)_{-1})=\phi(L(X)_{-1})=\phi(\mathfrak g_{-1})\subset L(X)_{-1}$,
and since~$L(X)$ is generated by $L(X)_{-1}$, we see that
$D(L(X)_{p})\subset L(X)_{p}$ and $D(\mathfrak a)\subset \mathfrak a$.
Thus there is a~derivation of $D_\phi$ of $\mathfrak m$ such that
$\pi\circ D=D_\phi\circ \pi$, where $\pi$ is the natural projection of $L(X)$ onto $\mathfrak m$.
The correspondence
$\gl(\mathfrak g_{-1})\ni\phi\mapsto D_\phi\in\Der(\mathfrak m)_0$
is an injective linear mapping.
Hence $\dim \gl(\mathfrak g_{-1})\leqq\dim \Der(\mathfrak m)_0$.
On the other hand, since $\mathfrak m$ is generated by $\mathfrak g_{-1}$,
the mapping $\Phi$ is a~Lie algebra monomorphism.
Therefore $\Phi$ is a Lie algebra isomorphism.
\end{proof}

\begin{remark}
Let $n$ and $\mu$ be positive integers with $n,\mu\geqq2$, and
let $\mathfrak m=\bigoplus\limits_{p<0}\mathfrak g_p$ be a free FGLA of
type $(n,\mu)$.  Furthermore
let $\mathfrak m'=\bigoplus\limits_{p<0}\mathfrak g'_p$
be an FGLA of the $\mu$-th kind,
and let $\varphi$ be a linear mapping of $\mathfrak g_{-1}$ into
$\mathfrak g'_{-1}$.
\begin{enumerate}\itemsep=0pt
\item
From the proof of Proposition \ref{prop2.1}, there exists a unique GLA
homomorphism $L(\varphi)$ of $\mathfrak m$ into $\mathfrak m'$
such that $L(\varphi)|\mathfrak g_{-1}=\varphi$.
\item
Let $\mathfrak m''=\bigoplus\limits_{p<0}\mathfrak g''_p$
be an FGLA of the $\mu$-th kind,
and let $\varphi'$ be a linear mapping of $\mathfrak g'_{-1}$ into~$\mathfrak g''_{-1}$. Assume that $\mathfrak m'=\bigoplus\limits_{p<0}\mathfrak g'_p$ is a free FGLA. By the uniqueness of $L(\varphi'\circ\varphi)$,
we see that
$L(\varphi'\circ\varphi)=L(\varphi')\circ L(\varphi)$.
\item Assume that $\mathfrak m'=\bigoplus\limits_{p<0}\mathfrak g'_p$ is a
free FGLA and $\varphi$ is injective. By the result of (2),
$L(\varphi)$ is a~monomorphism.
\item Let $W$ be an $m$-dimensional subspace of $\mathfrak g_{-1}$ with $m\geqq2$.
By the result of (3), the subalgebra of $\mathfrak m$ generated by $W$ is a
free FGLA of type $(m,\mu)$.
\end{enumerate}
\end{remark}
By Remark 2.1 (4) and \cite[Chapter~II, \S~2, Theorem~1]{bou72:1}, we get the following lemma.
\begin{lemma}\label{lem2.1} Let $\mathfrak m=\bigoplus\limits_{p<0}\mathfrak g_p$ be a
free FGLA of type $(n,\mu)$ with $\mu\geqq3$.
If $X$, $Y$ are linearly independent elements of $\mathfrak g_{-1}$, then
\begin{gather*}
\ad(X)^\mu(Y)=0, \qquad \ad(X)^{\mu-1}(Y)\ne 0, \\
\ad(Y)\ad(X)^{\mu-1}(Y)=0, \qquad \ad(Y)\ad(X)^{\mu-2}(Y)\ne 0.
\end{gather*}
\end{lemma}

\section{Universal fundamental graded Lie algebras}\label{section3}

Following N.~Tanaka \cite{tan70:1}, we introduce universal FGLAs of the
$\mu$-th kind.

Let $V$ be an $n$-dimensional vector space.
We def\/ine vector spaces $b(V)_p$ $(p<0)$ and li\-near mappings $B_p$ of
$\sum\limits_{r+s=p}b(V)_r\wedge b(V)_s$ into $b(V)_p$ $(p\leqq -2)$
as follows:
First of all, we put \mbox{$b(V)_{-1}=V$} and $b(V)_{-2}=\Lambda^2V$. Further
we def\/ine a mapping $B_{-2}:b(V)_{-1}\wedge b(V)_{-1}\to b(V)_{-2}$
to be the identity mapping. For $k\leqq-3$, we def\/ine $b(V)_k$ and
$B_k$ inductively as follows:
We set $b(V)^{(k+1)}=\bigoplus\limits_{p=-1}^{k+1}b(V)_p$ and
we def\/ine a subspace $c(V)_k$ of $\Lambda^2(b(V)^{(k+1)})$ to be
$\sum\limits_{r+s=k}b(V)_r\wedge b(V)_s$.
We denote by $A(V)_k$ the subspace of $c(V)_k$ spanned by the elements
\[
\underset{(X,Y,Z)}{\mathfrak S}\sum_{r+s=k}\sum_{u+v=r}B_r(X_u\wedge Y_v)\wedge Z_s,
\qquad  X,Y,Z\in b(V)^{(k+1)} ,
\]
where $\underset{(X,Y,Z)}{\mathfrak S}$ stands for the cyclic sum with respect to $X$, $Y$, $Z$, and $X_u$ denotes the $b(V)_u$-component in the decomposition
$b(V)^{(k+1)}=\bigoplus\limits_{p=-1}^{k+1}b(V)_p$.
Now we def\/ine $b(V)_k$ to be the factor space $c(V)_k/A(V)_k$, and
$B_k$ to be the projection of $c(V)_k$ onto $b(V)_k$.
We put $b(V)=\bigoplus\limits_{p<0}b(V)_p$ and def\/ine a bracket operation
$[\ , \ ]$ on $b(V)$ by
\[
[X,Y]=\sum_{p\leqq-2}\sum_{r+s=p}B_p(X_r\wedge Y_s)
\]
for all $X,Y\in b(V)$. Then $b(V)=\bigoplus\limits_{p<0}b(V)_p$ becomes a GLA generated by $b(V)_{-1}$, and $b(V)_p\ne0$ for all $p<0$ if $\dim V>1$.

Note that $b(V)_{-3}$ is isomorphic to $\Lambda^2(V)\otimes V/\Lambda^3V$.
Let $\mu$ be a positive integer. Assume that $\mu\geqq 2$ and $\dim V=n\geqq 2$.
Since $\bigoplus\limits_{p<-\mu}b(V)_p$ is a graded ideal of $b(V)$,
we see that the factor space $b(V,\mu)=b(V)/\bigoplus\limits_{p<-\mu}b(V)_p$
becomes an FGLA of $\mu$-th kind,
which is called a universal fundamental graded Lie algebra of the
$\mu$-th kind.
By \cite[Proposition~3.2]{tan70:1},
$b(V,\mu)$ is a free FGLA of type $(n,\mu)$.

\section{The prolongations of fundamental graded Lie algebras}\label{section4}

Following N.~Tanaka \cite{tan70:1}, we introduce the prolongations of FGLAs.
Let $\mathfrak m=\bigoplus\limits_{p<0}\mathfrak g_p$ be an FGLA.
A GLA
$\mathfrak g(\mathfrak m)
=\bigoplus\limits_{p\in\mathbb Z}\mathfrak g(\mathfrak m)_p$ is called the prolongation of $\mathfrak m$ if the following conditions hold:
\begin{enumerate}\itemsep=0pt
\renewcommand{\labelenumi}{$(\roman{enumi})$}
\item $\mathfrak g(\mathfrak m)_p=\mathfrak g_p$ for all $p<0$;
\item $\mathfrak g(\mathfrak m)$ is a transitive GLA;
\item If $\mathfrak h=\bigoplus\limits_{p\in\mathbb Z}\mathfrak h_p$ is a
GLA satisfying conditions (i) and (ii) above, then
$\mathfrak h=\bigoplus\limits_{p\in\mathbb Z}\mathfrak h_p$ can be embedded in
$\mathfrak g(\mathfrak m)$ as a GLA.
\end{enumerate}
We construct the prolongation $\mathfrak g(\mathfrak m)
=\bigoplus\limits_{p\in\mathbb Z}\mathfrak g(\mathfrak m)_p$ of $\mathfrak m$.
We set $\mathfrak g(\mathfrak m)_p=\mathfrak g_p$ $(p<0)$.
We def\/ine subspaces $\mathfrak g(\mathfrak m)_k$ $(k\geqq0)$ of
$\Hom(\mathfrak m,\bigoplus\limits_{p\leqq k-1}\mathfrak g(\mathfrak m)_p)_k$
and a bracket operation on $\mathfrak g(\mathfrak m)
=\bigoplus\limits_{p\in\mathbb Z}\mathfrak g(\mathfrak m)_p$ inductively.
First $\mathfrak g(\mathfrak m)_0$ is def\/ined to be $\Der(\mathfrak m)_0$ and a bracket operation $[\ ,\ ]:\bigoplus\limits_{p\leqq 0}\mathfrak g(\mathfrak m)_p\times \bigoplus\limits_{p\leqq 0}\mathfrak g(\mathfrak m)_p\to
\bigoplus\limits_{p\leqq 0}\mathfrak g(\mathfrak m)_p$ is def\/ined by
\begin{gather*}
[X,Y]=-[Y,X]=X(Y),\qquad  X\in\mathfrak g(\mathfrak m)_0, \quad Y\in\mathfrak m,\\
[X,Y]=XY-YX, \qquad  X,Y\in\mathfrak g(\mathfrak m)_0 .
\end{gather*}
Next for $k>0$ we def\/ine $\mathfrak g(\mathfrak m)_k$ $(k\geqq1)$ inductively
as follows:
\begin{gather*}
\mathfrak g(\mathfrak m)_k
=\Big\{ X\in\Hom\Big(\mathfrak m,\!\bigoplus_{p\leqq k-1}\mathfrak g(\mathfrak m)_p\Big)_k\!:
X([u,v])=[X(u),v]+[u,X(v)]\  \text{for all}\ u,v\in\mathfrak m\Big\},
\end{gather*}
where for $X\in\mathfrak g(\mathfrak m)_r$, $u\in\mathfrak m$,
we set $[X,u]=-[u,X]=X(u)$. Further for $X\in\mathfrak g(\mathfrak m)_k$,
$Y\in\mathfrak g(\mathfrak m)_l$ $(k,l\geqq0)$, by induction on $k+l\geqq0$,
we def\/ine $[X,Y]\in\Hom(\mathfrak m,\mathfrak g(\mathfrak m))_{k+l}$ by
\[
[X,Y](u)=[X,[Y,u]]-[Y,[X,u]],\qquad  u\in \mathfrak m .
\]
It follows easily that $[X,Y]\in \mathfrak g(\mathfrak m)_{k+l}$.
With this bracket operation, $\mathfrak g(\mathfrak m)
=\bigoplus\limits_{p\in\mathbb Z}\mathfrak g(\mathfrak m)_p$ becomes a graded
Lie algebra satisfying conditions $(i)$, $(ii)$ and $(iii)$ above.

Let $\mathfrak m$ and $\mathfrak g(\mathfrak m)$ be as above.
Assume that we are given a subalgebra $\mathfrak g_0$ of
$\mathfrak g(\mathfrak m)_0$.
We def\/ine subspaces $\mathfrak g_k$ $(k\geqq 1)$ of
$\mathfrak g(\mathfrak m)_k$ inductively as follows:
\[
\mathfrak g_k=\{ X\in\mathfrak g(\mathfrak m)_k:
[X,\mathfrak g_p]\subset \mathfrak g_{p+k}\
\text{for all}\ p<0 \}.
\]
If we put $\gla g$, then
$\gla g$ becomes a transitive graded Lie subalgebra of $\mathfrak g(\mathfrak m)$,
which is called the prolongation of
$(\mathfrak m,\mathfrak g_0)$.

By Proposition \ref{prop2.1} (2) we get the following proposition.
\begin{proposition}\label{prop4.1} Let $\mathfrak m=\bigoplus\limits_{p<0}\mathfrak g_p$ be a
free FGLA and let $\mathfrak g(\mathfrak m)
=\bigoplus\limits_{p\in\mathbb Z}\mathfrak g(\mathfrak m)_p$
be the prolongation of $\mathfrak m$.
Then the mapping $\mathfrak g(\mathfrak m)_0\ni D\mapsto D|\mathfrak g_{-1}
\in\gl(\mathfrak g_{-1})$ is an isomorphism.
\end{proposition}
Conversely we obtain the following proposition.

\begin{proposition}\label{prop4.2}
Let $\mathfrak m=\bigoplus\limits_{p<0}\mathfrak g_p$ be an
FGLA of the $\mu$-th kind and let
$\mathfrak g(\mathfrak m)
=\bigoplus\limits_{p\in\mathbb Z}\mathfrak g(\mathfrak m)_p$
be the prolongation of $\mathfrak m$.
Assume that $\mathfrak g(\mathfrak m)_0$ is isomorphic to
$\gl(\mathfrak g_{-1})$.
If $\mu=2$ or $\mu=3$,
then $\mathfrak m$ is a~free FGLA.
\end{proposition}

\begin{proof}\looseness=-1
We put $n=\dim \mathfrak g_{-1}$. We consider a universal FGLA
$b(\mathfrak g_{-1},\mu)=\bigoplus\limits_{p<0}b(\mathfrak g_{-1},\mu)_p$ of
the $\mu$-th kind. Since $b(\mathfrak g_{-1},\mu)$
is a free FGLA of type $(n,\mu)$,
there exists a GLA epimorphism $\varphi$
of $b(\mathfrak g_{-1},\mu)$ onto $\mathfrak m$ such that the restriction
$\varphi|b(\mathfrak g_{-1},\mu)_{-1}$ is the identity mapping.
Let $\check{b}(\mathfrak g_{-1},\mu)
=\bigoplus\limits_{p\in\mathbb Z}\check{b}(\mathfrak g_{-1},\mu)_p$
be the prolongation of $b(\mathfrak g_{-1},\mu)$.
Since the mapping $\mathfrak g(\mathfrak m)_0\ni D\mapsto
D|\mathfrak g_{-1}\in\gl(\mathfrak g_{-1})$ is an isomorphism,
$\varphi$ can be extended to be a homomorphism $\check{\varphi}$
of $\bigoplus\limits_{p\leqq0}\check{b}(\mathfrak g_{-1},\mu)_p$ onto
$\bigoplus\limits_{p\leqq0}\mathfrak g(\mathfrak m)_p$.
Let $\mathfrak a$ be the kernel of $\check{\varphi}$;
then $\mathfrak a$ is a graded ideal of $\bigoplus\limits_{p\leqq0}\check{b}(\mathfrak g_{-1},\mu)_p$. We set $\mathfrak a_p
=\mathfrak a\cap \check{b}(\mathfrak g_{-1},\mu)_p$; then
$\mathfrak a=\bigoplus\limits_{p\leqq 0}\mathfrak a_p$.
Since the restriction of $\check{\varphi}$ to $\check{b}(\mathfrak g_{-1},\mu)_{-1}
\oplus \check{b}(\mathfrak g_{-1},\mu)_0$ is injective,
$\mathfrak a_p=\{0\}$ for $p\geqq-1$. Also
each $\mathfrak a_p$ is a $\check{b}(\mathfrak g_{-1},\mu)_0$-submodule of
$\check{b}(\mathfrak g_{-1},\mu)_p$.
From the construction of $b(\mathfrak g_{-1},\mu)$, we see that
$b(\mathfrak g_{-1},\mu)_{-2}$ (resp. $b(\mathfrak g_{-1},\mu)_{-3}$) is
isomorphic to
$\Lambda^2(\mathfrak g_{-1})$
(resp.\ $\Lambda^2(\mathfrak g_{-1})\otimes \mathfrak g_{-1}/
\Lambda^3(\mathfrak g_{-1}))$ as a $\check{b}(\mathfrak g_{-1},\mu)_0$-module.
By the table of \cite{ov90:1},
$\Lambda^2(\mathfrak g_{-1})$ and $\Lambda^2(\mathfrak g_{-1})\otimes
\mathfrak g_{-1}/\Lambda^3(\mathfrak g_{-1})$ are irreducible
$\gl(\mathfrak g_{-1})$-modules.
Thus we see that $\mathfrak a_{-2}=\mathfrak a_{-3}=\{0\}$.
From $\mu\leqq3$ it follows that $\varphi$ is an isomorphism.
\end{proof}


\section{Finite-dimensional simple graded Lie algebras}\label{section5}

Following \cite{yam93:1}, we f\/irst state the classif\/ication of
f\/inite-dimensional simple GLAs.

\looseness=-1
Let $\gla g$ be a f\/inite-dimensional simple GLA of the $\mu$-th kind over
$\mathbb C$ such that the negative part $\mathfrak g_-$ is an FGLA.
Let $\mathfrak h$ be a Cartan subalgebra of $\mathfrak g_0$;
then $\mathfrak h$ is a Cartan subalgebra of $\mathfrak g$ such that
$E\in \mathfrak h$, where
$E$ is  the element of $\mathfrak g_0$ such that
$[E,x]=px$ for all $x\in\mathfrak g_p$ and~$p$.
Let~$\Delta$ be
a root system of $(\mathfrak g,\mathfrak h)$. For~$\alpha\in\Delta$,
we denote by $\mathfrak g^{\alpha}$ the root space corresponding to~$\alpha$. We set
$\mathfrak h_\mathbb R=\{ h\in\mathfrak h:\alpha(h)\in\mathbb R\ \text{for all}\ \alpha\in \Delta \}$
and let $(h_1,\dots,h_l)$ be a basis of $\mathfrak h_\mathbb R$ such that
$h_1=E$.
We def\/ine the set of positive roots $\Delta^+$ as the set of roots
which are positive with respect to the lexicographical ordering
in $\mathfrak h_{\mathbb R}^*$ determined by the basis $(h_1,\dots,h_l)$ of
$\mathfrak h_{\mathbb R}$.
Let $\Pi\subset \Delta^+$ be the corresponding simple root system.
We denote by $\{m_1,\dots,m_l\}$ the coordinate functions
corresponding to $\Pi$, i.e., for $\alpha\in\Delta$,
we can write $\alpha=\sum\limits_{i=1}^l m_i(\alpha)\alpha_i$.

We set $\alpha_i(E)=s_i$ and $\bm s=(s_1,\dots,s_l)$;
then each $s_i$ is a non-negative integer.
For $\alpha\in\Delta$, we call the integer
$\ell_{\bm s}(\alpha)=\sum\limits_{i=1}^l m_i(\alpha)s_i$ the
$\bm s$-length of $\alpha$.
We put $\Delta_p=\{ \alpha\in\Delta: \ell_{\bm s}(\alpha)=p \}$,
$\Pi_p=\Delta_p\cap \Pi$ and $I=\{ i\in\{1,\dots,l\}: s_i=1 \}$.
Let $\theta$ be the highest root of $\mathfrak g$; then
$\ell_{\bm s}(\theta)=\mu$. Also since the $\mathfrak g_0$-module
$\mathfrak g_{-\mu}$ is irreducible,
$\dim \mathfrak g_{-\mu}=1$
if and only if $\langle \theta,\alpha\spcheck_i\rangle=0$ for all
$i\in\{1,\dots,l\}\setminus I$, where $\{\alpha\spcheck_i\}$ is the simple root system of the dual root system $\Delta\spcheck$ of $\Delta$ corresponding
to $\{\alpha_i\}$.
In our situation, since $\mathfrak g_-$ is generated by
$\mathfrak g_{-1}$, we have $s_i=0$ or 1 for all $i$.
The $l$-tuple $\bm s=(s_1,\dots,s_l)$ of non-negative integers
is determined only by the ordering of $(\alpha_1,\dots,\alpha_l)$.
In what follows, we assume that the ordering of
$(\alpha_1,\dots,\alpha_l)$ is as in the table of~\cite{bou68:1}.
If $\mathfrak g$ has the Dynkin diagram of type $X_l$ $(X=A,\dots,G)$,
then the simple GLA $\gla g$
is said to be of type $(X_l,\Pi_1)$.
Here we remark that for an automorphism $\bar{\mu}$ of the
Dynkin diagram, a~simple GLA of type $(X_l,\Pi_1)$
is isomorphic to that of type $(X_l,\bar{\mu}(\Pi_1))$.
We will identify a~simple GLA of type $(X_l,\Pi_1)$
with that of type $(X_l,\bar{\mu}(\Pi_1))$.

For $i\in I$, we put $\Delta_p^{(i)}=\{\alpha\in\Delta:
m_i(\alpha)=p\ \text{and}\ m_j(\alpha)=
0\ \text{for}\ j\in I\setminus\{i\} \}$ and
$\mathfrak g_p^{(i)}=\sum\limits_{\alpha\in\Delta_p^{(i)}}\mathfrak g^\alpha$;
then
$\mathfrak g_{-1}^{(i)}$ is an irreducible
$\mathfrak g_0$-submodule of $\mathfrak g_{-1}$ with
highest weight $-\alpha_i$.
In particular, if the $\mathfrak g_0$-module $\mathfrak g_{-1}$
is irreducible, then $\#(I)=1$.

{\samepage For $i\in I$,
we denote by $\mathfrak g^{(i)}$ the subalgebra of $\mathfrak g$ generated by
$\mathfrak g_{-1}^{(i)}\oplus \mathfrak g_{1}^{(i)}$;
then $\mathfrak g^{(i)}$ is a~simple GLA whose
Dynkin diagram is the connected component containing the vertex $i$ of
the sub\-diagram of~$X_l$
corresponding to vertices $(\{1,\dots,l\}\setminus I)\cup\{i\}$.
We denote by $\theta^{(i)}$ the highest root of $\mathfrak g^{(i)}$.
Then $[\mathfrak g_{-1}^{(i)},\mathfrak g_{-1}^{(i)}]=\{0\}$ if and only if
$m_i(\theta^{(i)})=1$.

}

From Theorem 5.2 of \cite{yam93:1}, we obtain the following theorem:
\begin{theorem}\label{thm5.1}
Let $\gla g$ be a finite-dimensional simple GLA over $\mathbb C$ such that
$\mathfrak g_-$ is an FGLA and
the $\mathfrak g_0$-module $\mathfrak g_{-1}$
is irreducible.
Then $\gla g$ is the prolongation of $\mathfrak g_-$ except for
the following cases:
\begin{enumerate}\itemsep=0pt
\renewcommand{\labelenumi}{$(\alph{enumi})$}
\item $\mathfrak g_-$ is of the first kind;
\item $\mathfrak g_-$ is of the second kind and $\dim\mathfrak g_{-2}=1$.
\end{enumerate}
\end{theorem}

Let $\gla g$ be a f\/inite-dimensional simple GLA.
Now we assume that $\mathfrak g_0$ is isomorphic to $\gl(\mathfrak g_{-1})$;
then the $\mathfrak g_0$-module $\mathfrak g_{-1}$ is irreducible.
The derived subalgebra $[\mathfrak g_0,\mathfrak g_0]$ of $\mathfrak g_0$
is a~semisimple Lie algebra whose Dynkin diagram is the subdiagram of
$X_l$ consisting
of the vertices $\{1,\dots,l\}\setminus I$.
Since $[\mathfrak g_0,\mathfrak g_0]$ is of type $A_{l-1}$ and since
the $\mathfrak g_0$-module $\mathfrak g_{-1}$ is elementary,
$(X_l,\Delta_1)$ is one of the following cases:
\[
(A_l,\{\alpha_1\}),\qquad (B_l,\{\alpha_l\}), \quad  l\geqq2,
\qquad (G_2,\{\alpha_1\}).
\]
From this result and Propositions \ref{prop4.1} and \ref{prop4.2}, we get the following theorem:
\begin{theorem}\label{thm5.2}
Let $\gla g$ be a finite-dimensional simple GLA of type
$(X_l,\Pi_1)$ over $\mathbb C$ satisfying the following conditions:
\begin{enumerate}\itemsep=0pt
\renewcommand{\labelenumi}{$(\roman{enumi})$}
\item $\mathfrak g_-$ is an FGLA of the $\mu$-th kind;
\item The $\mathfrak g_0$-module $\mathfrak g_{-1}$ is irreducible;
\item $\mathfrak g_0$ is isomorphic to $\gl(\mathfrak g_{-1})$;
\item $\mathfrak g$ is the prolongation of $\mathfrak g_{-}$.
\end{enumerate}

Then $\mathfrak g_-$ is a free FGLA of type
$(l,\mu)$, and
$\gla g$ is one of the following types:
\begin{enumerate}\itemsep=0pt
\renewcommand{\labelenumi}{$(\alph{enumi})$}
\item $l\geqq3$, $\mu=2$, $(X_l,\Pi_1)=(B_l,\{\alpha_l\})$.
\item $l=2$, $\mu=3$, $(X_l,\Pi_1)=(G_2,\{\alpha_1\})$.
\end{enumerate}
\end{theorem}

\section[Graded Lie algebras $W(n)$, $K(n)$ of Cartan type]{Graded Lie algebras $\boldsymbol{W(n)}$, $\boldsymbol{K(n)}$ of Cartan type}\label{section6}

In this section, following V.G.~Kac \cite{kac68:1},
we describe Lie algebras $W(n)$, $K(n)$ of Cartan type and their standard gradations.

Let $A(m)$ denote the monoid (under addition) of all $m$-tuples of
non-negative integers.
For an $m$-tuple
$\bm s=(s_1,\dots,s_m)$ of positive integers and
$\alpha=(\alpha_1,\dots,\alpha_m)\in A(m)$
we set $\|\alpha\|_{\bm s}=\sum\limits_{i=1}^ms_i\alpha_i$.
Also we denote the $m$-tuple $(1,\dots,1)$ by ${\bm 1}_m$
and  we denote the
$(m+1)$-tuple $(1,\dots,1,2)$ by $({\bm 1}_m,2)$.
Let $\mathfrak A(m)=\mathbb C[x_1,\dots,x_m]$.
For any $m$-tuple $\bm s$
of positive integers, we denote by $\mathfrak A(m;\bm s)_p$
the subspace of $\mathfrak A(m)$ spanned by polynomials
\[
x^{\alpha}=x_1^{\alpha_1}\cdots x_m^{\alpha_m}, \qquad
 \alpha=(\alpha_1,\dots,\alpha_m)\in A(m),\quad  \|\alpha\|_{\bm s}=p .
\]
Let $W(m)$ be the Lie algebra consisting of all the polynomial vector f\/ields
\begin{gather}\label{eq6.1}
\sum\limits_{i=1}^m P_i\frac{\partial}{\partial x_i},
\qquad  P_i\in\mathfrak A(m) .
\end{gather}

For an $m$-tuple $\bm s=(s_1,\dots,s_m)$ of positive integers,
we denote by $W(m;\bm s)_p$
the subspaces of $W(m)$ consisting of those polynomial vector f\/ields \eqref{eq6.1}
such that the polynomials $P_i$ are contained in
$\mathfrak A(m;\bm s)_{p+s_i}$;
then
$W(m;\bm s)=\bigoplus\limits_{p\in\mathbb Z}W(m;\bm s)_p$
is a transitive GLA.
In particular, $W(m;{\bm 1}_m)=\bigoplus\limits_{p\geqq-1}W(m;{\bm 1}_m)_p$ is
a transitive irreducible GLA
such that: $(i)$ $W(m;{\bm 1}_m)_0$ is isomorphic to $\gl(m,\mathbb C)$;
$(ii)$ the $W(m;{\bm 1}_m)_0$-module $W(m;{\bm 1}_m)_{-1}$ is elementary; $(iii)$
$W(m;{\bm 1}_m)$ is the prolongation of $W(m;{\bm 1}_m)_-$.

We now consider the following dif\/ferential form
\[
\omega_K=dx_{2n+1}-\sum_{i=1}^nx_{i+n}dx_{i}.
\]
Def\/ine
\[
K(n)=\{ D\in W(2n+1):D\omega_K\in\mathfrak A(2n+1)\omega_K \}.
\]
(Here the action of $D$ on the dif\/ferential forms is extended from
its action $\mathfrak A(2n+1)$ by requiring that $D$ be derivation of
the exterior algebra satisfying $D(df)=d(Df)$, where
$df=\sum\frac{\partial f}{\partial x_i}dx_i$, $f\in\mathfrak A(m)$.)
We set $K(n)_p=W(2n+1;({\bm 1}_{2n},2))_p\cap K(n)$.
Then $K(n)=\bigoplus\limits_{p\geqq-2}K(n)_p$
is a transitive irreducible GLA
such that: $(i)$
$K(n)_0$ is isomorphic to $\csp(n,\mathbb C)$;
$(ii)$ the $K(n)_0$-module $K(n)_{-1}$ is elementary;
$(iii)$ $K(n)$ is the prolongation of $K(n)_-$
(cf.~\cite{kac68:1, mor88:1}).

From Proposition 2.2 of \cite{mt70:1}, we get
\begin{theorem}\label{thm6.1} Let $\gla g$ be a transitive GLA over $\mathbb C$
satisfying the following conditions:
\begin{enumerate}\itemsep=0pt
\renewcommand{\labelenumi}{$(\roman{enumi})$}
\item $\mathfrak g_-$ is an FGLA of the
$\mu$-th kind;
\item $\mathfrak g$ is infinite-dimensional;
\item The $\mathfrak g_0$-module $\mathfrak g_{-1}$ is irreducible;
\item $\mathfrak g$ is the prolongation of $\mathfrak g_-$.
\end{enumerate}
Then $\mu\leqq2$ and $\gla g$ is isomorphic to $W(m;{\bm 1}_m)$ or $K(n)$.
\end{theorem}
